\chapter{Introduction} \label{Chap1}
\pagenumbering{arabic}

\section{Motivation}

Methane (CH\textsubscript{4}) is a powerful greenhouse gas, with a global warming potential
roughly 28 times that of CO\textsubscript{2} over a 100-year horizon and around 80 times that of
CO\textsubscript{2} over a 20-year horizon \cite{ipcc2021}. Agriculture is the leading sectoral
source of anthropogenic methane, accounting for approximately 49\% of agricultural greenhouse gas
emissions through enteric fermentation alone in the UK and a comparably large share globally
\cite{fao2023, eea2024}, and global livestock emissions are projected to keep rising alongside
demand for animal products \cite{fao2023}. Predicting CH\textsubscript{4} flux at the ecosystem
scale is genuinely difficult: emissions are driven by the continuous interaction of multiple,
non-stationary biophysical factors -- soil moisture, temperature, rainfall -- together with
discrete livestock management interventions that dominate the signal in UK managed grassland
specifically.

The North Wyke Farm Platform (NWFP) offers a rare and underutilised resource to address this
difficulty: continuous half-hourly Eddy Covariance (EC) flux measurements from managed beef
cattle grassland spanning over seven years. To the author's knowledge, no prior study has applied
ML-based forecasting to EC CH\textsubscript{4} flux at a managed temperate grassland -- existing
work in this space addresses gap-filling (within-distribution interpolation) rather than
forecasting (sequential extrapolation), a structurally different modelling problem
(Chapter~\ref{Chap2}). Closing this gap matters beyond academic interest: accurate prediction and
scenario projection can support farm-level management decisions -- grazing timing, stocking
density -- and inform scenario and policy planning toward national emissions-reduction targets.

This dissertation addresses the gap through a three-stage pipeline: gap-filling a continuous
CH\textsubscript{4} series from NWFP's EC record (Chapter~\ref{Chap4}), forecasting that series
using a benchmarked, uncertainty-quantified model roster (Chapter~\ref{Chap5}), and using the
resulting models to project CH\textsubscript{4} response under contrasting management and climate
scenarios (Chapter~\ref{Chap6}).

\section{Research Questions}

\begin{table}[h]
\centering
\begin{tabular}{cp{11cm}}
\toprule
ID & Question \\
\midrule
RQ1 & What ML approaches exist for CH\textsubscript{4} flux prediction in agricultural/grassland
contexts; do any address ecosystem-scale EC data; how does gap-filling differ from multi-step
forecasting as a modelling constraint? \\
RQ2 & How do statistical baselines, tree-based ensembles, and deep learning compare for
half-hourly EC CH\textsubscript{4} forecasting under temporal variability and non-stationarity? \\
RQ3 & How transferable are findings from adjacent domains (wetland gap-filling, animal-scale
prediction) to managed grassland EC forecasting? \\
RQ4 & How can XAI (SHAP) and UQ (quantile ML / conformal prediction) decode interactions between
continuous environmental variables and discrete management interventions? \\
RQ5 & What structural requirements are needed to integrate forecasting models into a digital
shadow with ``what-if'' scenario analysis at farm scale? \\
\bottomrule
\end{tabular}
\caption{Research questions guiding this dissertation.}
\end{table}

\section{Aims and Objectives}

The overall aim of this dissertation is the development and evaluation of multiple machine
learning approaches for CH\textsubscript{4} forecasting at NWFP, and to demonstrate their
integration with scenario projection under contrasting management and climate conditions. This
aim is pursued through five objectives:

\begin{enumerate}
    \item Conduct a systematic literature review on time-series forecasting, quantile machine
    learning, eddy covariance sensor methodology, and agricultural digital twin frameworks.
    \item Acquire, pre-process, and document NWFP EC CH\textsubscript{4} data (2018--present),
    including gap imputation, quality control, and feature engineering.
    \item Design and implement a benchmarking pipeline comparing statistical baselines, tree
    ensembles, and deep learning architectures, using temporal cross-validation to prevent data
    leakage.
    \item Apply interpretability methods (SHAP) and uncertainty quantification (conformal
    prediction / quantile machine learning) to characterise model behaviour.
    \item Generate synthetic management and climate scenarios and evaluate model behaviour under
    them.
\end{enumerate}

\section{Contributions}

This dissertation contributes three benchmarked results to the field:

\begin{enumerate}
    \item \textbf{A gap-filling benchmark for EC CH\textsubscript{4} flux at a managed
    grassland} (Chapter~\ref{Chap4}), comparing established methods (MDS, Random Forest),
    published baseline replications, and newer tabular/foundation model families, with
    uncertainty quantification and interpretability attached. The best configuration reaches an
    hourly-resolution $R^2_{OLS}$ of 0.41--0.60 across all three towers (Zhu et al.'s own reporting
    convention, matched directly for comparability), against a published ceiling of $R^2<0.1$ at
    hourly resolution for the same ecosystem type \cite{zhu2023a}.

    \item \textbf{A recursive multi-step forecasting benchmark for EC CH\textsubscript{4} flux}
    (Chapter~\ref{Chap5}), the first, to the author's knowledge, applied to this ecosystem-scale
    signal at a managed temperate grassland, comparing statistical baselines, tree ensembles, and
    tabular foundation models under identical conditions, with calibrated uncertainty
    quantification. The best configuration -- an equal-weight blend of three TabPFN variants,
    including a conservative spike-excess correction -- beats a day-of-year climatology baseline by
    a wide margin ($\text{MASE}=0.691$), though a dedicated spike-vs-non-spike evaluation shows this
    advantage is concentrated on ordinary days and essentially disappears on the high-magnitude
    events that dominate the CH\textsubscript{4} budget (Chapter~\ref{Chap5}). Interpretability
    work across both gap-filling and forecasting
    converges on a single, repeatedly-confirmed finding: livestock density is the dominant driver
    of CH\textsubscript{4} flux at NWFP, consistent across SHAP attribution, native feature
    importance, SARIMAX coefficients, and TabPFN permutation importance.

    \item \textbf{Scenario projection under contrasting management and climate conditions}
    (Chapter~\ref{Chap6}): the forecasting architecture developed in this dissertation is
    extended to produce distinguishable, interpretable CH\textsubscript{4} projections under
    synthetic livestock and climate scenarios, grounded in CMIP6-based climate data and
    externally-anchored, literature-plausible stocking levels.
\end{enumerate}

\section{Report Overview}

Chapter~\ref{Chap2} reviews the literature on EC flux gap-filling, forecasting, agricultural
digital twins, and recent developments in tabular modelling, establishing the gap this
dissertation addresses. Chapter~\ref{Chap3} describes the NWFP dataset and the preprocessing
pipeline built on top of it. Chapter~\ref{Chap4} presents the gap-filling benchmark that produces
a continuous CH\textsubscript{4} series for training and evaluation, including its uncertainty
quantification. Chapter~\ref{Chap5} presents the recursive forecasting benchmark and its
uncertainty quantification, this dissertation's primary modelling contribution. Chapter~\ref{Chap6}
presents scenario projection under contrasting management and climate conditions.
Chapter~\ref{Chap7} discusses the findings and their limitations, and Chapter~\ref{Chap8}
concludes.
